\abstract{
	\IfFileExists{Abstract_Section.md.tex}{
		\input{Abstract_Section.md.tex} % This loads the file. This is relative to the location of this file.
	}
	{} % If the file *doesn't exist, do nothing.
	This dissertation presents two dark sector searches with the ATLAS experiment using proton-proton collision collected at the Large Hadron Collider (LHC), one with the full Run 2 dataset and one with the first 134.6 fb\(^{-1}\) of Run 3 data.
	The first search is for dark mesons decaying into top and bottom quarks ($tttb$ or $ttbb$), which is sensitive to a proposed Stealth Dark Matter scenario with a strongly coupled dark sector and a confining gauge group.
	This model contains a viable dark matter candidate in the form of a scalar baryon. 
	The search targets both the all-hadronic and one-lepton final states of the top quark decays, with the one-lepton channel providing the strongest limits to date on this model.
	Dark pions with masses below 940 GeV are excluded for a ratio between the dark pion mass and the dark rho mass of 0.45, and masses below 740 GeV for a ratio of 0.25.
	The second search is for the exotic decay of the Higgs boson into a photon and a dark photon, which is motivated by several BSM models and provides an additional opportunity to probe a hypothetical dark sector.
	The analysis is optimized for the gluon-gluon fusion production channel, which has an order of magnitude larger cross-section than the vector boson fusion and associated production channels, and is made possible by the implementation of a new composite trigger that greatly increases the acceptance of gluon-gluon fusion events.
	The results of this search set the most stringent limits on the branching ratio of Higgs to a photon and dark photon to date, setting a combined upper limit on the branching ratio better than 1\%. 
	The search for BSM physics in dark sectors will continue to be a major focus of the ATLAS physics program in the coming years, and the results of these searches will be used to guide future searches for dark matter at the LHC and beyond.
	
	This dissertation includes previously published coauthored material.
}
